% !TEX root = ../algorithms.tex

\section{Универсальные семейства хеш-функций. 2-универсальность семейства multiply-shift. Хеш-таблицы с разрешением коллизий цепочками}

\begin{Definition}{Универсальное и $c$-универсальное семейство хеш-функций}{}
	Семейство хеш-функций $\mathcal{H}$ над универсумом $U$, где каждая $h \in \mathcal{H}$ имеет вид $h\colon U \to \{0,\ldots,m-1\}$, называется \emph{универсальным}, если
	\[
	\forall x,y \in U,\ x \neq y\colon\qquad
	\Pr_{h \in \mathcal{H}}\bigl(h(x) = h(y)\bigr) \leq \frac{1}{m},
	\]
	где вероятность берётся по случайному равновероятному выбору $h \in \mathcal{H}$.
	
	Семейство $\mathcal{H}$ называется \emph{$c$-универсальным} (для $c \geq 1$), если
	\[
	\forall x,y \in U,\ x \neq y\colon\qquad
	\Pr_{h \in \mathcal{H}}\bigl(h(x) = h(y)\bigr) \leq \frac{c}{m}.
	\]
\end{Definition}

\begin{Remark}{}{}
	Чтобы выбрать $h$ случайно и равновероятно из множества \emph{всех} функций вида $U \to \{0,\ldots,m-1\}$, потребовалось бы заполнить таблицу размера $|U|$ --- это нереалистично. Идея состоит в том, чтобы выбирать $h$ случайно и равновероятно из небольшого семейства $\mathcal{H} \subseteq \{0,\ldots,m-1\}^{U}$ с малой вероятностью коллизий.
\end{Remark}

\subsection{Семейство multiply-shift}

Пусть ключи --- $w$-битовые целые числа: $U = \{0,1,\ldots,2^w-1\}$, $m = 2^t$.

\begin{Definition}{Семейство multiply-shift}{}
	\[
	\mathcal{H}_{\mathrm{ms}} = \{h_a : a \in U,\ a \text{ нечётно}\},
	\qquad
	h_a(x) = \bigl(a \cdot x \bmod 2^w\bigr) \gg (w-t).
	\]
	Чтобы выбрать $h_a \in \mathcal{H}_{\mathrm{ms}}$ случайно и равновероятно:
	\begin{enumerate}
		\item положить $a = \mathrm{random}(0,\ldots,2^w-1)$;
		\item положить $a := a \mathbin{|} 1$ (побитовое \textsc{или} с $1$ делает $a$ нечётным).
	\end{enumerate}
\end{Definition}

\begin{Remark}{}{}
	Другое распространённое семейство --- \emph{multiply-add} (линейное хеширование по простому модулю): для произвольного $U = \{0,\ldots,u-1\}$ фиксируем простое $p \geq u$ и полагаем
	\[
	h_{a,b}(x) = \bigl((ax + b) \bmod p\bigr) \bmod m,
	\qquad
	\mathcal{H}_{\mathrm{ma}} = \{h_{a,b} : a,b \in \{0,\ldots,p-1\},\ a \neq 0\}.
	\]
\end{Remark}

\begin{Theorem}{2-универсальность семейства multiply-shift}{}
	Семейство $\mathcal{H}_{\mathrm{ms}} = \{h_a\}$ является $2$-универсальным.
\end{Theorem}

\begin{proof}
	Пусть $x,y \in U = \{0,1,\ldots,2^w-1\}$, $x < y$, и $m = 2^t$. По определению $h_a$,
	\[
	h_a(x) = h_a(y)
	\quad\Longleftrightarrow\quad
	ax \bmod 2^w \text{ и } ay \bmod 2^w \text{ имеют одинаковые старшие } t \text{ бит}.
	\]
	Равенство старших $t$ бит чисел $ax \bmod 2^w$ и $ay \bmod 2^w$ означает, что старшие $t$ бит числа $a(y-x) \bmod 2^w$ равны либо $\underbrace{00\ldots0}_{t}$, либо $\underbrace{11\ldots1}_{t}$ (второй случай возникает из-за возможного переноса при вычитании).
	
	Запишем $y - x = z \cdot 2^b$, где $z$ нечётно и $0 \leq b \leq w-1$. Рассмотрим два случая.
	
	\textbf{Случай $1^\circ$: $b > w - t$.}
	Так как $a$ и $z$ нечётны, число $az$ нечётно, и
	\[
	a(y-x) \bmod 2^w = \bigl(az \bmod 2^{w-b}\bigr) \cdot 2^b
	\]
	является нечётным кратным $2^b$. Поэтому в двоичной записи $a(y-x) \bmod 2^w$ бит номер $b$ равен $1$, а биты номеров $0,\ldots,b-1$ равны $0$. Поскольку $b > w-t$, бит $b$ принадлежит старшим $t$ битам (номера $w-t,\ldots,w-1$), и среди тех же старших $t$ бит найдётся бит номера $b-1 \geq w-t$, равный $0$. Значит, старшие $t$ бит не могут быть равны ни $00\ldots0$ (есть единица), ни $11\ldots1$ (есть ноль) при любом нечётном $a$, то есть $\Pr_a\bigl(h_a(x) = h_a(y)\bigr) = 0$.
	
	\textbf{Случай $2^\circ$: $b \leq w - t$.}
	Запишем нечётное $a$ в виде $a = 2a_1 + 1$, где $a_1$ --- случайное равновероятное $(w-1)$-битовое число (это эквивалентно случайному равновероятному выбору нечётного $a \in \{0,\ldots,2^w-1\}$). Тогда
	\[
	a(y-x) = (2a_1+1) z \cdot 2^b = (z + 2a_1 z) \cdot 2^b.
	\]
	Отображение $a_1 \mapsto a_1 z \bmod 2^{w-1}$ является биекцией на множестве $(w-1)$-битовых чисел: если $a_1 < a_2$ --- два различных $(w-1)$-битовых числа, то $a_2 - a_1 \not\equiv 0 \pmod{2^{w-1}}$, а так как $z$ нечётно (и значит обратимо по модулю $2^{w-1}$), то $z(a_2-a_1) \not\equiv 0 \pmod{2^{w-1}}$, то есть $a_1 z \not\equiv a_2 z \pmod{2^{w-1}}$.
	
	Следовательно, при случайном равновероятном выборе $a_1$ величина $a_1 z \bmod 2^{w-1}$ также случайна и равновероятна на множестве $(w-1)$-битовых чисел, а значит младшие $(w-1)$ бит числа $a(y-x) \bmod 2^w$ распределены случайно и равновероятно. Поэтому старшие $t$ бит числа $a(y-x) \bmod 2^w$ равны $\underbrace{00\ldots0}_{t}$ с вероятностью не более $\tfrac{1}{2^t}$ и равны $\underbrace{11\ldots1}_{t}$ с вероятностью не более $\tfrac{1}{2^t}$. Суммируя,
	\[
	\Pr_a\bigl(h_a(x) = h_a(y)\bigr) \leq \frac{2}{2^t} = \frac{2}{m}.
	\]
	В обоих случаях $\Pr_a\bigl(h_a(x) = h_a(y)\bigr) \leq \tfrac{2}{m}$, то есть семейство multiply-shift является $2$-универсальным.
\end{proof}

\subsection{Хеш-таблицы с разрешением коллизий цепочками}

\begin{Definition}{}{}
	Пусть задана таблица $T[0,\ldots,m-1]$ и хеш-функция $h\colon U \to \{0,\ldots,m-1\}$. В хеш-таблице с разрешением коллизий цепочками каждая ячейка $T[k]$ хранит связный список всех ключей $x$, для которых $h(x) = k$. Операции $\mathrm{insert}$, $\mathrm{delete}$, $\mathrm{query}$ с ключом $x$ обращаются к ячейке $T[h(x)]$ и просматривают (или обновляют) её список.
\end{Definition}

\begin{figure}[H]
	\centering
	\begin{tikzpicture}[>=Stealth, font=\small, cell/.style={draw,minimum width=1cm,minimum height=0.8cm}]
		\node[cell] (t0) at (0,0) {$0$};
		\node[cell] (t1) at (0,-0.8) {$1$};
		\node (td) at (0,-1.6) {$\vdots$};
		\node[cell] (tm) at (0,-2.4) {$m-1$};
		\node[above] at (0,0.6) {$T$};
		
		\node[cell] (a1) at (2.2,0) {$x_1$};
		\node[cell] (a2) at (3.6,0) {$x_2$};
		\draw[->] (t0.east) -- (a1.west);
		\draw[->] (a1.east) -- (a2.west);
		\draw[->] (a2.east) -- ++(0.7,0) node[right] {$\varnothing$};
		
		\draw[->] (t1.east) -- ++(0.7,0) node[right] {$\varnothing$};
		
		\node[cell] (b1) at (2.2,-2.4) {$y_1$};
		\draw[->] (tm.east) -- (b1.west);
		\draw[->] (b1.east) -- ++(0.7,0) node[right] {$\varnothing$};
	\end{tikzpicture}
	\caption{Хеш-таблица с цепочками: ячейка $T[k]$ хранит связный список всех ключей $x$ с $h(x)=k$. Время операции с ключом $x$ определяется длиной списка в ячейке $T[h(x)]$.}
\end{figure}

\begin{Statement}{Перевыделение таблицы}{}
	Зафиксируем константу $\alpha > 1$ (например, $\alpha = 2$). Если после очередной вставки число ключей $n$ в таблице размера $m$ превышает $\alpha m$, выделяется новая таблица размера $m' = \beta m$ при некотором $\beta > 1$ (например, $\beta = 2$); все ключи переносятся в новую таблицу, после чего старая таблица удаляется. Эта стратегия гарантирует, что в любой момент $n \leq \alpha m$.
\end{Statement}

\begin{Theorem}{}{}
	Пусть $h$ выбирается случайно и равновероятно из $c$-универсального семейства $\mathcal{H}$, где $c = O(1)$, а размер таблицы поддерживается стратегией перевыделения (так что $n \leq \alpha m$). Тогда ожидаемое амортизированное время каждой из операций равно $O(1)$.
\end{Theorem}

\begin{proof}
	\textbf{Часть 1. Ожидаемое время операции без учёта перестройки.}
	Достаточно показать, что для произвольного $x \in U$ ожидаемая длина списка $T[h(x)]$ равна $O(1)$.
	
	Пусть $y_1,\ldots,y_n$ --- ключи, хранящиеся в таблице. Для каждого $i$ положим
	\[
	I_{y_i} =
	\begin{cases}
		0, & h(y_i) \neq h(x),\\
		1, & h(y_i) = h(x).
	\end{cases}
	\]
	Тогда длина списка $T[h(x)]$ равна $\sum_{i=1}^{n} I_{y_i}$, и по линейности математического ожидания и $c$-универсальности:
	\begin{align*}
		\mathbf{E}\!\left[\sum_{i=1}^{n} I_{y_i}\right]
		= \sum_{i=1}^{n} \Pr_{h \in \mathcal{H}}\bigl(h(x) = h(y_i)\bigr)
		\leq n \cdot \frac{c}{m}
		\leq (\alpha m) \cdot \frac{c}{m}
		= \alpha c = O(1).
	\end{align*}
	
	\textbf{Часть 2. Амортизированная стоимость перестройки.}
	Перестройка происходит, когда $n > \alpha m$: создаётся таблица размера $\beta m$ ($\beta > 1$ фиксировано), все $n \leq \alpha m + 1$ ключей перехешируются. Стоимость одной перестройки --- $O(m)$.
	
	Покажем, что между двумя последовательными перестройками выполняется не менее $\Omega(m)$ операций вставки. Пусть сразу после перестройки к таблице размера $\beta m$ в ней хранится $n_0$ ключей. Так как перестройка была инициирована условием $n > \alpha m$, до неё размер был $m$, значит $n_0 \leq \alpha m + 1$. Следующая перестройка произойдёт лишь когда число ключей превысит $\alpha \cdot \beta m$, то есть потребуется не менее
	\[
	\alpha \beta m - n_0 \;\geq\; \alpha \beta m - \alpha m - 1 \;=\; \alpha(\beta - 1)\,m - 1
	\]
	операций вставки. Стоимость перестройки $O(\beta m)$, распределяя её по этим вставкам:
	\[
	\frac{O(\beta m)}{\alpha(\beta-1)\,m} = O\!\left(\frac{\beta}{\alpha(\beta-1)}\right) = O(1),
	\]
	поскольку $\alpha$ и $\beta$ --- фиксированные константы, $\beta > 1$.
\end{proof}